\section*{Оценка европейских опционов в шестипериодической биномиальной модели}

\subsection*{Исходные параметры}

Рассмотрим базовый актив со следующими параметрами:
\begin{itemize}
    \item Начальная цена: $S_0 = 100$
    \item Число периодов: $N = 6$
    \item Множитель роста: $u = 1.2$
    \item Множитель падения: $d = 0.8$
    \item Безрисковая ставка: $r = 5\%$
\end{itemize}

\subsection*{Риск-нейтральная вероятность}

В биномиальной модели риск-нейтральная вероятность определяется как:

:contentReference[oaicite:0]{index=0}

Подставим значения:

\[
p = \frac{1.05 - 0.8}{1.2 - 0.8} = \frac{0.25}{0.4} = 0.625
\]

\subsection*{Цена актива в конечных узлах}

Цена актива в узле с $k$ подъёмами:

\[
S_T(k) = S_0 \cdot u^k \cdot d^{N-k}
\]

\subsection*{Выплаты опционов}

\begin{itemize}
    \item Call:
    \[
    C_T = \max(S_T - K, 0)
    \]
    \item Put:
    \[
    P_T = \max(K - S_T, 0)
    \]
\end{itemize}

\subsection*{Цена опциона}

Цена опциона определяется как математическое ожидание дисконтированных выплат:

:contentReference[oaicite:1]{index=1}

Аналогично для put:

\[
P_0 = \frac{1}{(1+r)^N} \sum_{k=0}^{N} \binom{N}{k} p^k (1-p)^{N-k} \max(K - S_T(k), 0)
\]

\subsection*{Выводы}

\begin{itemize}
    \item Цена call возрастает с уменьшением страйка $K$
    \item Цена put возрастает с увеличением страйка
    \item При $K \approx S_0$ опционы имеют наибольшую чувствительность
    \item Модель даёт приближение к рыночным ценам, но:
    \begin{itemize}
        \item не учитывает волатильность как непрерывный процесс
        \item чувствительна к выбору $u$ и $d$
    \end{itemize}
\end{itemize}
